\section*{Задание}

Создать таблицу для сравнения генерации псевдослучайных чисел табличным и алгоритмическим способами.
Таблица должна содержать три столбца: числа, сгенерированные табличным способом, числа, сгенерированные алгоритмическим способом, произвольные числа, введенные пользователем (10 чисел).
Разработать количественный критерий случайности для оценки последовательностей.
Для каждого из трех столбцов вычислить и отобразить значение этого критерия, которое будет численно характеризовать степень "случайности" соответствующей последовательности.

\newpage

\section*{Теоретическая часть}

Критерий "Лунные фазы" оценивает случайность последовательности чисел на основе их связи с текущей фазой Луны. Метод определяет текущую фазу Луны по дню месяца и проверяет, насколько распределение чисел, кратных номеру этой фазы, соответствует статистическим ожиданиям.

\subsection*{Алгоритм работы}

\begin{enumerate}
\item Определение фазы Луны (на основе дня месяца вычисляется текущая фаза Луны с присвоением номера от 1 до 8 и учётом реальной длительности каждой фазы).

\item Подсчёт кратных чисел.

\item Взвешивание по редкости фазы (более короткие фазы (например, Новолуние - 1 день) получают больший весовой коэффициент, так как обнаружение кратных чисел в эти периоды статистически менее вероятно).

\item Сравнение фактического соотношения кратных чисел с ожидаемым соотношением, рассчитанным на основе длительности фазы. Учитывается только избыток кратных чисел сверх ожидаемого уровня.

\end{enumerate}

\subsection*{Основная формула оценки}

\[
S(N,D) = 1 - \min\left(1, \dfrac{\max\left(0, \dfrac{C(N,p)}{|N|} - \dfrac{d}{L}\right) \cdot \min\left(3, \dfrac{L}{d}\right)}{1 - \dfrac{d}{L}}\right)
\]

где:
\begin{itemize}
\item $C(N,p)$ -- количество чисел, кратных номеру фазы
\item $\dfrac{d}{L}$ -- ожидаемая доля (идеальное соотношение)
\item $1 - \dfrac{d}{L}$ -- максимально возможный избыток
\end{itemize}


\subsection*{Область значений}

$S = 1$ -- максимальная случайность (нет избытка кратных чисел).
$S = 0$ -- минимальная случайность (максимальный избыток кратных чисел).


Формула учитывает как статистическое отклонение, так и редкость лунной фазы через весовой коэффициент.


\section*{Результат}

\begin{figure}[H]
\centering
\includegraphics[width=0.97\textwidth]{img/table.jpg}
\caption{Интерфейс программы}
\end{figure}